% =========================================================
% The Complex Numbers
% =========================================================

\section{The Complex Numbers}

\begin{definition}[Complex Numbers]
\label{def:complex-numbers}
The set of \emph{complex numbers}, denoted $\mathbb{C}$, is defined as
\[
\mathbb{C} := \{(a,b) : a,b \in \mathbb{R}\},
\]
with operations defined by
\[
(a,b) + (c,d) := (a+c,\; b+d),
\]
\[
(a,b)\cdot(c,d) := (ac - bd,\; ad + bc).
\]
\end{definition}

\begin{remark*}[Fully quantified form]
\[ \mathbb{C} := \{a+bi : a,b\in\mathbb{R}\},\; i^2=-1. \]
\end{remark*}

\begin{remark*}[Flash]
\FlashQ{Define the complex numbers.}
\FlashA{$\mathbb{C}=\{a+bi:a,b\in\mathbb{R}\}$ where $i^2=-1$.}
\end{remark*}

\begin{remark*}
We identify $(a,b)$ with the symbol
\[
a + bi,
\]
where $i := (0,1)$ satisfies $i^2 = -1$.
\end{remark*}

\begin{remark*}[Field structure of $\mathbb{C}$]
With the operations defined above, $\mathbb{C}$ satisfies the field axioms:

\begin{itemize}
\item $(\mathbb{C}, +)$ is an abelian group,
\item $(\mathbb{C}\setminus\{0\}, \cdot)$ is an abelian group,
\item multiplication distributes over addition.
\end{itemize}

Hence $\mathbb{C}$ is a field containing $\mathbb{R}$ as the subset
$\{(a,0): a \in \mathbb{R}\}$.
\end{remark*}